\chapter{Mar.~10 --- Sheaf Cohomology}

\section{Motivation for Sheaf Cohomology}
\begin{theorem}\label{thm:coh}
  If $\mathcal{G}$ is a coherent sheaf
  on a projective variety $X$, then
  $\Gamma(X, \mathcal{G})$ is a finite
  dimensional $k$-vector space.
\end{theorem}

\begin{remark}
  Theorem \ref{thm:coh} will be a kind of
  motivating theorem for sheaf cohomology.
  We will try to give a simple proof.
  To do this, we use the following strategy:
  \begin{enumerate}
    \item There is a closed embedding
      $i : X \hookrightarrow \PP^n$.
      Now $\Gamma(X, \mathcal{G}) = \Gamma(\PP^n, i_*\mathcal{G})$.
      Note that $i_* \mathcal{G}$ is
      coherent since $i$ is a closed
      embedding. So we reduce to the case
      $X = \PP^n$.
    \item When $\mathcal{G} = \bigoplus_{i = 1}^r \OO_{\PP^n}(a_i)$,
      Theorem \ref{thm:coh} holds as
      \[
        \Gamma\Big(\PP^n, \bigoplus_{i = 1}^r \OO_{\PP^n}(a_i)\Big)
        = \bigoplus_{i = 1}^r k[x_0, \dots, x_n]_{a_i}.
      \]
    \item We want to reduce to Case (2).
      By a homework problem, there is a
      surjection
      \[
        \begin{tikzcd}
          \mathcal{F}
          = \bigoplus_{i = 1}^r \OO_{\PP^n}(a_i)
          \ar[r, two heads] & \mathcal{G}
          \ar[r] & 0
        \end{tikzcd}
      \]
      for some $a_1, \dots, a_r$. Set
      $\mathcal{E} = \ker(\mathcal{F} \to \mathcal{G})$, then
      we get a short exact sequence
      \[
        \begin{tikzcd}
          0 \ar[r] & \mathcal{E} \ar[r] & \mathcal{F}
          \ar[r] & \mathcal{G} \ar[r] & 0
        \end{tikzcd}
      \]
      Now we get an exact sequence
      \[
        \begin{tikzcd}
          0 \ar[r] & \Gamma(\PP^n, \mathcal{E})
          \ar[r] & \Gamma(\PP^n, \mathcal{F})
          \ar[r] & \Gamma(\PP^n, \mathcal{G})
        \end{tikzcd}
      \]
      as taking global sections is left
      exact. Note that $\Gamma(\PP^n, \mathcal{F})$
      is finite dimensional. Though we do
      not always get surjectivity onto
      $\Gamma(\PP^n, \mathcal{G})$, we
      can construct functors
      \begin{align*}
        \QCoh_X \longrightarrow \mathrm{Vect}_k \\
        \mathcal{F} \longmapsto H^i(X, \mathcal{F})
      \end{align*}
      for $i \ge 0$ such that:
      \begin{enumerate}[(i)]
        \item $H^0(X, \mathcal{F}) \cong \Gamma(X, \mathcal{F})$,
        \item $H^i(X, \mathcal{F}) = 0$ for
          $i > \dim X$,
        \item given a short exact sequence
          $0 \to \mathcal{E} \to \mathcal{F} \to \mathcal{G} \to 0$
          in $\QCoh_X$, there is a long
          exact sequence
          \[
            \begin{tikzcd}
              H^i(X, \mathcal{E}) \ar[r] & H^i(X, \mathcal{F})
              \ar[r] & H^i(X, \mathcal{G})
              \ar[r] & H^{i + 1}(X, \mathcal{E})
              \ar[r] & \cdots
            \end{tikzcd}
          \]
        \item $H^i(\PP^n, \bigoplus_{j = 1}^r \OO_{\PP^n}(a_j))$
          is a finite dimensional $k$-vector space.
      \end{enumerate}
      Then to prove the theorem, it suffices
      to show that
      for any coherent sheaf $\mathcal{G}$
      on $\PP^n$ and $i \ge 0$,
      \[H^i(\PP^n, \mathcal{G})\]
      is a finite
      dimensional $k$-vector space.
      The theorem follows from this by
      property (i).

      Assuming we have $H^i$ satisfying
      (i)-(iv), we can prove the above
      claim via downwards induction.
      As the base case, note that for
      $i > n$, we have
      $H^i(\PP^n, \mathcal{G}) = 0$ for
      all quasicoherent sheaves $\mathcal{G}$ on $\PP^n$ by (ii).
      For the inductive step, assume
      the claim holds for all coherent
      $\mathcal{G}$
      and $i > i_0$. Consider
      \[
        \begin{tikzcd}
          0 \ar[r] & \mathcal{E} \ar[r] & \mathcal{F} = \bigoplus_{j = 1}^r \OO_{\PP^n}(a_j)
          \ar[r] & \mathcal{G} \ar[r] & 0
        \end{tikzcd}
      \]
      and the corresponding long exact
      sequence (by (iii))
      \[
        \begin{tikzcd}
          H^{i_0}(\PP^n, \mathcal{E}) \ar[r] & H^{i_0}(\PP^n, \mathcal{F})
          \ar[r] & H^{i_0}(\PP^n, \mathcal{G})
          \ar[r] & H^{i_0 + 1}(\PP^n, \mathcal{E})
            \ar[r] & \cdots
        \end{tikzcd}
      \]
      Note that $H^{i_0}(X, \mathcal{F})$
      is finite dimensional by (iv)
      and $H^{i_0 + 1}(X, \mathcal{E})$ is finite dimensional
      by the inductive hypothesis.
      So
      $H^{i_0}(X, \mathcal{G})$ is finite dimensional, and by induction
      we get the full result.
  \end{enumerate}
  Thus our goal is to give a definition
  of $H^i$, which we will do
  using the Cech complex.
\end{remark}

\section{Homological Algebra}

\begin{definition}
  Let $R$ be a ring. A \emph{cochain complex}
  of $R$-modules is a sequence of
  $R$-modules
  \[
    A = A^\bullet
    := (A^0 \overset{\partial}{\longrightarrow} A^1 \overset{\partial}{\longrightarrow} A^2 \overset{\partial}{\longrightarrow} \cdots)
  \]
  with morphisms
  $\partial^i : A^i \to A^{i + 1}$ such
  that $\partial^{i + 1} \circ \partial^i = 0$.
\end{definition}

\begin{definition}
  Let $A^\bullet$ be a cochain complex.
  The \emph{cohomology} of $A^\bullet$ is
  \[
    H^i(A^\bullet)
    = \frac{\ker \partial^i}{\im \partial^{i - 1}}.
  \]
  We call $Z^i(A^\bullet) = \ker \partial^i$ the
  \emph{cocycles} and $B^i(A^\bullet) = \im \partial^{i - 1}$ the
  \emph{coboundaries}.
\end{definition}

\begin{definition}
  A \emph{morphism of cochain complexes}
  $A^\bullet \to B^\bullet$ is the data of
  morphisms $\phi^i : A^i \to B^i$ such that
  the following diagram commutes for
  all $i$:
  \[
    \begin{tikzcd}
      A^i \ar[r, "\partial"] \ar[d, swap, "\phi^i"] & A^{i + 1} \ar[d, "\phi^{i + 1}"] \\
      B^i \ar[r, "\partial"] & B^{i + 1}
    \end{tikzcd}
  \]
\end{definition}

\begin{remark}
  Given a morphism of cochain
  complexes $A^\bullet \to B^\bullet$,
  there are induced maps
  \[
    Z^i(A^\bullet) \to Z^i(B^\bullet)
    \quad \text{and} \quad
    B^i(A^\bullet) \to B^i(B^\bullet).
  \]
  In particular, we get an induced
  map on cohomology
  $H^i(A^\bullet) \to H^i(B^\bullet)$.
\end{remark}

\begin{definition}
  A \emph{short exact sequence of cochain complexes}
  \[
    \begin{tikzcd}
      0 \ar[r] & A^\bullet \ar[r, "\phi"] & B^\bullet \ar[r, "\psi"] & C^\bullet \ar[r] & 0
    \end{tikzcd}
  \]
  is given by morphisms $\phi$ and $\psi$
  such that
  \[
    \begin{tikzcd}
      0 \ar[r] & A^i \ar[r, "\phi^i"] & B^i \ar[r, "\psi^i"] & C^i \ar[r] & 0
    \end{tikzcd}
  \]
  is exact for all $i \ge 0$.
\end{definition}

\begin{remark}
  The category of cochain complexes
  of $R$-modules is an abelian category.
\end{remark}

\begin{prop}\label{prop:long}
  Given a short exact sequence of
  cochain complexes
  $0 \to A^\bullet \to B^\bullet \to C^\bullet \to 0$, we get a long exact
  sequence
  in cohomology
  \[
    \begin{tikzcd}
      H^i(A^\bullet) \ar[r] & H^i(B^\bullet) \ar[r] & H^i(C^\bullet) \ar[r] & H^{i + 1}(A^\bullet) \ar[r] & \cdots
    \end{tikzcd}
  \]
\end{prop}

\begin{proof}
  The difficulty is to define
  the map $H^{i - 1}(C^\bullet) \to H^i(A^\bullet)$.
  Consider the diagram
  \[
    \begin{tikzcd}
      & H^{i - 1}(A) \ar[r] \ar[d] & H^{i - 1}(B) \ar[r] \ar[d] & H^{i - 1}(C) \ar[d] \\
      & A^i / B^i(A) \ar[r] \ar[d, swap, "\partial"] & B^i / B^i(B) \ar[r] \ar[d, swap, "\partial"] & C^i / B^i(C) \ar[r] \ar[d, swap, "\partial"] & 0 \\
      0 \ar[r] & Z^i(A) \ar[r] \ar[d] & Z^i(B)
      \ar[r] \ar[d] & Z^i(C) \ar[d] \\
             & H^i(A) \ar[r] & H^i(B) \ar[r] & H^i(C)
    \end{tikzcd}
  \]
  The snake lemma then gives a map
  $H^{i - 1}(C) \to H^i(A)$.
\end{proof}

\section{Cech Cohomology}

\begin{definition}
  Let $X$ be a topological space,
  $\mathcal{F}$ a sheaf of abelian groups
  on $X$, and $\mathcal{U} = \{U_i\}_{i = 0}^n$
  an open cover of $X$.
  For $I \subseteq \{0, \dots, n\}$, denote
  $U_I = \bigcap_{i \in I} U_i$. For
  $p \ge 0$, set
  \[
    C^p(\mathcal{U}, X)
    = \prod_{\substack{I \subseteq \{0, \dots, n\} \\ |I| = p + 1}}
    \mathcal{F}(U_I).
  \]
  Then the \emph{Cech complex} is the
  cochain complex
  \[
    C^\bullet(\mathcal{U}, \mathcal{F})
    = [
    C^0(\mathcal{U}, \mathcal{F}) \overset{d^0}{\longrightarrow} C^1(\mathcal{U}, \mathcal{F}) \overset{d^1}{\longrightarrow} \cdots
    ]
  \]
  with $d^p : C^p(\mathcal{U}, \mathcal{F}) \to C^{p + 1}(\mathcal{U}, \mathcal{F})$ given by
  $(s_I)_I \mapsto (s_J)_J$, where
  $J = \{i_0 < i_1 < \cdots < i_{p + 1}\}$
  and
  \[
    s_J = \sum (-1)^k s_{J \setminus \{i_k\}}|_{U_J}.
  \]
\end{definition}

\begin{exercise}
  Check that the Cech complex is in
  fact a cochain complex, i.e. that
  $d^{i + 1} \circ d^i = 0$.
\end{exercise}

\begin{definition}
  The \emph{Cech cohomology} is
  given by $H^i(\mathcal{U}, \mathcal{F}) = H^i(C^\bullet(\mathcal{U}, \mathcal{F}))$.
\end{definition}

\begin{example}
  We compute $H^0(\mathcal{U}, \mathcal{F})$.
  We have
  \begin{align*}
    H^0(\mathcal{U}, \mathcal{F})
    &= \ker(C^0(\mathcal{U}, \mathcal{F}) \to C^1(\mathcal{U}, \mathcal{F})) \\
    &= \ker\Big(
      \prod_{i \in I} \mathcal{F}(U_i)
      \to \prod_{0 \le i < j \le n} \mathcal{F}(U_{i, j})
    \Big)
    \cong \mathcal{F}(X)
    = \Gamma(X, \mathcal{F})
  \end{align*}
  where the isomorphism is by the sheaf
  axiom.
\end{example}

\begin{example}
  Let $X = \PP^1$ and
  $\mathcal{U} = (U_0, U_1)$ be the
  standard open cover.
  \begin{enumerate}[(i)]
    \item Let $\mathcal{F} = \OO_X$. Then
      we have
      \begin{align*}
        C^\bullet(\mathcal{U}, \OO_X)
        = [\OO_{\PP^1}(U_0) \oplus \OO_{\PP^1}(U_1)
        \to \OO_{\PP_{1}}(U_0 \cap U_1)]
        = [k[y / x] \oplus k[x / y] \to k[x / y, y / x]],
      \end{align*}
      where the map $\alpha : k[y / x] \oplus k[x / y] \to k[x / y, y / x]$ is given by
      $(f, g) \mapsto f - g$. Thus
      \[
        H^0(\mathcal{U}, \OO_{\PP^1})
        = \ker \alpha \cong k \cong \Gamma(\PP^1, \OO_{\PP^1}).
      \]
      As the map $\alpha$ is surjective,
      we have that
      \[
        H^1(\mathcal{U}, \OO_{\PP^1})
        = \coker \alpha = 0.
      \]
    \item Let $\mathcal{F} = \Omega_{\PP^1}$.
      Then we have the cochain complex
      \begin{align*}
        C^\bullet(\mathcal{U}, \Omega_{\PP^1})
        &= [\Omega_{\PP^1}(U_0) \oplus \Omega_{\PP^1}(U_1) \to \Omega_{\PP^1}(U_{0, 1})] \\
        &= [k[y / x] d(y / x) \oplus k[x / y] d(x / y) \to d(y / x) k[x / y, y / x]].
      \end{align*}
      If we set $s = y / x$, then the above
      becomes
      \[
        C^\bullet(\mathcal{U}, \OO_{\PP^1})
        = [k[s] ds \oplus k[s^{-1}] ds^{-1} \to k[s^{\pm 1}] ds],
      \]
      where the map $\beta : k[s] ds \oplus k[s^{-1}] ds^{-1} \to k[s^{\pm 1}] ds$ is given by
      \[
        (f(s) ds, g(s^{-1}) ds^{-1})
        \longmapsto f(s) ds - g(s^{-1})\bigg({-\frac{1}{s^2}}\bigg) ds.
      \]
      Thus we see that
      \begin{align*}
        H^0(\mathcal{U}, \Omega_{\PP^1})
        &= \ker \beta = 0 \\
        H^1(\mathcal{U}, \Omega_{\PP^1})
        &= \coker \beta
        = \Span\{\overline{(1 / s) ds}\}
        \cong k.
      \end{align*}
  \end{enumerate}
\end{example}

\begin{example}
  Let $X$ be a variety and
  $\mathcal{F} = \OO_X^*$ (this is a
  sheaf of groups under multiplication). Fix an open cover
  $\mathcal{U} = \{U_i\}_{i = 0}^n$ of $X$.
  Consider
  \[
    H^1(\mathcal{U}, \OO_X^*)
    = \{
      (g_{i, j})_{i, j}
      : g_{i, j} \in \OO_X^*(U_{i, j})
      \text{ and } g_{i, j} g_{j, k} g_{i, k}^{-1} = 1
    \} / {\sim},
  \]
  where $(g_{i, j})_{i, j} \sim (g'_{i, j})_{i, j}$
  if there exists $h_i \in \OO_X^*(U_i)$
  such that $h_i^{-1} g_{i, j} h_j = g'_{i, j}$. Thus
  \[
    H^1(\mathcal{U}, \OO_X^*)
    \cong \{\text{line bundles } \mathbb{L} \to X \text{ such that } \mathbb{L}_{U_i} \cong \mathbbm{1}_{U_i}\} / {\cong}.
  \]
  In fact, the above
  is an isomorphism of groups.
\end{example}

\begin{remark}
  From now on, we will assume that
  $X$ is a variety, $\mathcal{F}$
  is a quasicoherent sheaf on $X$, and that
  $\mathcal{U} = \{U_i\}_{i = 0}^n$ is an
  open affine cover of $X$.
\end{remark}

\begin{prop}
  We have the following:
  \begin{enumerate}
    \item $H^i(\mathcal{U}, \mathcal{F}) = 0$ for $i > n$, where
      $n = |\mathcal{U}| - 1$.
    \item Given a short exact sequence
      $0 \to \mathcal{E} \to \mathcal{F} \to \mathcal{G} \to 0$
      in $\QCoh_X$, there is an induced
      long exact sequence
      on Cech cohomology:
      \[
        \begin{tikzcd}
          H^i(\mathcal{U}, \mathcal{E}) \ar[r] & H^i(\mathcal{U}, \mathcal{F})
          \ar[r] & H^i(\mathcal{U}, \mathcal{G})
          \ar[r] & H^{i + 1}(\mathcal{U}, \mathcal{E}) \ar[r] & \cdots
        \end{tikzcd}
      \]
  \end{enumerate}
\end{prop}

\begin{proof}
  (1) We have $C^p(\mathcal{U}, \mathcal{F}) = 0$ for $p > n$.
  
  (2) Since the $U_I$ are affine, we get a
  short exact sequence
  \[
    \begin{tikzcd}
      0 \ar[r] & \mathcal{E}(U_I) \ar[r] & \mathcal{F}(U_I) \ar[r] & \mathcal{G}(U_I) \ar[r] & 0
    \end{tikzcd}
  \]
  for any $I \subseteq \{0, \dots, n\}$.
  So we get an induced short exact sequence
  of Cech complexes
  \[
    \begin{tikzcd}
      0 \ar[r] & C^\bullet(\mathcal{U}, \mathcal{E}) \ar[r] & C^\bullet(\mathcal{U}, \mathcal{F}) \ar[r] & C^\bullet(\mathcal{U}, \mathcal{G}) \ar[r] & 0
    \end{tikzcd}
  \]
  This gives a long exact sequence on
  Cech cohomology by Proposition
  \ref{prop:long}.
\end{proof}
